\section{Motivation and main mechanisms}
The recurring strategy has three steps.  First, normalize the source expression until the object being tested is a kernel, moment sequence, or logarithmic derivative.  Second, prove positivity, complete monotonicity, or a finite determinant inequality in that normal form.  Third, transport the result back to the source problem, often obtaining a sharp endpoint, a monotonicity interval, or a counterexample.
This organization also explains the negative results.  When the positive-kernel route cannot hold, an endpoint expansion, a rational specialization, or an explicit numerical witness often identifies exactly where the proposed statement fails.
